\section{Fourier series}

\subsection{Orthogonal functions}

For real functions on $[a,b]$, define the inner product by
\begin{equation}
(f,g)=\int_a^b f(x)g(x)\,dx.
\end{equation}
The Schwarz inequality is
\begin{equation}
|(f,g)|^2\leq(f,f)(g,g),
\end{equation}
with equality precisely when $f$ and $g$ are linearly dependent. Functions are orthogonal when $(f,g)=0$; the squared norm is $\lVert f\rVert^2=(f,f)$. For complex functions, the appropriate inner product is
\begin{equation}
(f,g)=\int_a^b f^*(x)g(x)\,dx.
\end{equation}
A set $\{\varphi_n\}$ is orthonormal when $(\varphi_m,\varphi_n)=\delta_{mn}$. Orthogonality implies linear independence: multiplying $\sum_n c_n\varphi_n=0$ by $\varphi_j^*$ and integrating gives $c_j=0$.

Just as for vectors, a linearly independent function set can be orthogonalized. The Gram--Schmidt construction begins with $\varphi_1=f_1$ and continues as
\begin{equation}
\varphi_n=f_n-\sum_{j=1}^{n-1}
\frac{(f_n,\varphi_j)}{(\varphi_j,\varphi_j)}\varphi_j.
\end{equation}
For an orthonormal set, the best finite expansion $f\approx\sum_{n=1}^N c_n\varphi_n$ has coefficients $c_n=(f,\varphi_n)$. Positivity of the squared error gives Bessel's inequality,
\begin{equation}
\sum_{n=1}^{\infty}|c_n|^2\leq\lVert f\rVert^2.
\end{equation}
An orthonormal set is complete when its finite expansions approximate every piecewise continuous function arbitrarily well in mean square. For a complete set, Bessel's inequality becomes Parseval's relation,
\begin{equation}
\sum_{n=1}^{\infty}c_n^*d_n=(f,g),
\qquad c_n=(f,\varphi_n),\quad d_n=(g,\varphi_n).
\end{equation}

This is exactly the same idea as projecting a vector onto an orthonormal basis in Euclidean space: the Fourier coefficients are the components of the function in the chosen basis, and orthogonality makes the components independent.

\subsection{Origin in heat flow}

Fourier introduced trigonometric series while studying heat flow in a uniform rod of length $\ell$. In dimensionless units, a temperature $T(x,t)$ with zero endpoint temperatures obeys
\begin{equation}
T_{xx}=T_t,\qquad T(0,t)=T(\ell,t)=0.
\end{equation}
Letting $T=f(x)g(t)$ gives $f''/f=g'/g=-k^2$. The boundary conditions require $k=n\pi/\ell$ and yield the separated solutions
\begin{equation}
T_n(x,t)=A_n\sin\left(\frac{n\pi x}{\ell}\right)
\exp\left(-\frac{n^2\pi^2t}{\ell^2}\right).
\end{equation}
Linearity permits their superposition. Thus an arbitrary initial temperature $T(x,0)=f(x)$ is represented by a sine series,
\begin{equation}
f(x)=\sum_{n=1}^{\infty}A_n\sin\left(\frac{n\pi x}{\ell}\right),
\qquad
A_n=\frac{2}{\ell}\int_0^\ell f(x)\sin\left(\frac{n\pi x}{\ell}\right)\,dx.
\end{equation}

\subsection{Fourier theorem}

For a function of period $2\ell$, the Fourier series is
\begin{equation}
f(x)=\frac{a_0}{2}+\sum_{n=1}^{\infty}
\left[a_n\cos\left(\frac{n\pi x}{\ell}\right)+b_n\sin\left(\frac{n\pi x}{\ell}\right)\right],
\end{equation}
where
\begin{align}
a_n&=\frac{1}{\ell}\int_{-\ell}^{\ell}f(x)\cos\left(\frac{n\pi x}{\ell}\right)\,dx,\\
b_n&=\frac{1}{\ell}\int_{-\ell}^{\ell}f(x)\sin\left(\frac{n\pi x}{\ell}\right)\,dx.
\end{align}
The trigonometric basis is orthogonal on one period:
\begin{equation}
(\sin nx,\sin mx)=\pi\delta_{mn},\qquad
(\cos nx,\cos mx)=\pi\delta_{mn},\qquad
(\sin nx,\cos mx)=0
\end{equation}
for $m,n\neq0$. This orthogonality is the reason the coefficients can be extracted by simple integration. The Dirichlet conditions require finitely many discontinuities and extrema in each period. At a jump $x_0$, the series converges to $[f(x_0^-)+f(x_0^+)]/2$. Odd functions have sine series; even functions have cosine series.

The equivalent complex form is
\begin{equation}
f(x)=\sum_{n=-\infty}^{\infty}c_ne^{\mathrm{i}n\pi x/\ell},
\qquad c_n=\frac{1}{2\ell}\int_{-\ell}^{\ell}f(x)e^{-\mathrm{i}n\pi x/\ell}\,dx.
\end{equation}
For real $f$, $c_{-n}=c_n^*$. Parseval's identity becomes
\begin{equation}
\frac{1}{2\ell}\int_{-\ell}^{\ell}|f(x)|^2\,dx=\sum_{n=-\infty}^{\infty}|c_n|^2
=\frac{a_0^2}{4}+\frac12\sum_{n=1}^{\infty}(a_n^2+b_n^2).
\end{equation}

\begin{example}
On $(-\ell,\ell)$ let $f(x)=x$ and extend it periodically. This odd sawtooth has
\begin{equation}
f(x)=\frac{2\ell}{\pi}\sum_{n=1}^{\infty}
\frac{(-1)^{n+1}}{n}\sin\left(\frac{n\pi x}{\ell}\right).
\end{equation}
At $x=\ell/2$, this gives the Leibniz series $\pi/4=1-1/3+1/5-\cdots$. Applying Parseval's identity to a square wave similarly gives $\sum_{n=0}^{\infty}(2n+1)^{-2}=\pi^2/8$. The underlying idea is again the same: any sufficiently regular periodic function can be decomposed into a sum of elementary harmonics, and the coefficients control the weight of each harmonic.
\end{example}

